\documentclass[11pt]{article}
\newcommand{\ArticleShort}{Foundations}
\usepackage[T1]{fontenc}
\usepackage[utf8]{inputenc}
\usepackage[english]{babel}
\usepackage{newtxtext}
\usepackage{amsmath,amsthm,mathtools}
\usepackage{newtxmath}
\usepackage[letterpaper,margin=1in,headheight=15pt]{geometry}
\usepackage{microtype,setspace,booktabs,tabularx,longtable,array,enumitem,titlesec,fancyhdr}
\usepackage{xcolor}
\definecolor{ink}{HTML}{193244}
\usepackage{csquotes}
\usepackage[backend=biber,style=apa,natbib=true,doi=true,url=true,eprint=false]{biblatex}
\addbibresource{shared/references.bib}
\usepackage[unicode,hidelinks,bookmarksnumbered=true]{hyperref}
\usepackage{bookmark,xurl}
\setstretch{1.07}
\setlength{\parindent}{1.3em}
\setlength{\parskip}{0.2em}
\setlength{\emergencystretch}{2em}
\setlength{\bibitemsep}{0.45\baselineskip}
\setlength{\bibhang}{1.5em}
\renewcommand*{\bibfont}{\small}
\setcounter{biburllcpenalty}{7000}
\setcounter{biburlucpenalty}{7000}
\setcounter{biburlnumpenalty}{7000}
\setlist{topsep=0.4em,itemsep=0.25em,parsep=0pt}
\titleformat{\section}{\large\bfseries\color{ink}}{\thesection}{0.65em}{}
\titleformat{\subsection}{\normalsize\bfseries}{\thesubsection}{0.65em}{}
\titlespacing*{\section}{0pt}{1.2\baselineskip}{0.45\baselineskip}
\titlespacing*{\subsection}{0pt}{0.8\baselineskip}{0.3\baselineskip}
\pagestyle{fancy}
\fancyhf{}
\fancyhead[L]{\small\itshape \AE{}ther-flow and General Relativity}
\fancyhead[R]{\small\ArticleShort}
\fancyfoot[L]{\scriptsize Revised manuscript \textperiodcentered{} 9 September 2026}
\fancyfoot[R]{\small\thepage}
\renewcommand{\headrulewidth}{0.3pt}
\renewcommand{\footrulewidth}{0pt}
\newcommand{\Aether}{\AE{}ther}
\newcommand{\Aflow}{\AE{}ther-flow}
\newcommand{\TheoryName}{\Aflow{} Interpretation of Relativity}
\newcommand{\TheoryProgram}{\Aether{} / \Aflow{} framework}
\newcommand{\Sub}{\mathcal A}
\newcommand{\PhiSrc}{\Phi_{\mathrm{src}}}
\newcommand{\Order}{\mathcal S}
\newcommand{\Ph}{\Phi}
\newcommand{\Proj}{D}
\newcommand{\TT}{\mathrm{TT}}
\newcommand{\dd}{\mathrm d}
\newcommand{\R}{\mathbb R}
\newtheorem{definition}{Definition}
\newtheorem{proposition}{Proposition}
\theoremstyle{remark}
\newtheorem{remark}{Remark}
\newcommand{\PaperTitle}[3]{%
  \hypersetup{pdftitle={#1},pdfsubject={#2},pdfauthor={}}%
  \thispagestyle{plain}%
  \begin{center}
  {\small\scshape \AE{}ther-flow and General Relativity\par}
  \vspace{0.7em}
  {\LARGE\bfseries\color{ink}#1\par}
  \vspace{0.5em}
  \if\relax\detokenize{#2}\relax\else
  {\normalsize #2\par}
  \vspace{0.35em}
  \fi
  \end{center}
  \begin{abstract}\noindent #3\end{abstract}
  \vspace{0.35em}
}
\newcommand{\References}{\printbibliography[heading=bibintoc,title={References}]}

% Copyright footer added 10 September 2026.
\fancyfoot[L]{\scriptsize Revised manuscript \textperiodcentered{} 15 September 2026 \textperiodcentered{} Copyright \textcopyright{} 2026 Alexander Samuel Ricciardi \textperiodcentered{} AngryOwl}
\fancyfoot[R]{\small\thepage}
\fancypagestyle{plain}{%
  \fancyhf{}%
  \fancyfoot[L]{\scriptsize Revised manuscript \textperiodcentered{} 15 September 2026 \textperiodcentered{} Copyright \textcopyright{} 2026 Alexander Samuel Ricciardi \textperiodcentered{} AngryOwl}%
  \fancyfoot[R]{\small\thepage}%
  \renewcommand{\headrulewidth}{0pt}%
  \renewcommand{\footrulewidth}{0pt}%
}

\begin{document}
\PaperTitle{Ontology: Foundations}{}{The \Aether{} / \Aflow{} framework asks whether spatial and temporal relations can be recovered from intrinsic ordered motion in a proposed four-dimensional source structure, rather than assumed at the outset. Its present assumptions specify a smooth four-dimensional manifold but leave the mathematical type of the source flow unresolved. General relativity is adopted separately as the quantitative gravitational benchmark. This paper defines the source vocabulary, separates interpretive aims from mathematical results, and sets the conditions for assessing a future reconstruction. The research has two purposes: investigate a source explanation of relativity and evaluate whether an artificial-intelligence research system can produce independently checkable deductions or scoped negative results. A counterexample to entailment, an obstruction within a specified model class, and an unresolved search have different meanings. No source-to-spacetime derivation or general impossibility theorem is established here.}

\section{The question and the present answer}
The \TheoryProgram{} is motivated by a question about physical explanation: could the spacetime relations described by relativity arise from a more basic organization, rather than being the starting point of the account? The words \emph{\Aether{}} and \emph{\Aflow{}} express that aim. They do not, by themselves, specify the objects or laws needed to answer it.

The project uses this source question as both a physics investigation and a proposed AI-assisted research case study. Its physical aim is a deeper account of GR, not a departure from its predictions introduced for its own sake. GR agreement is the recovery target against which source candidates are tested, not a guaranteed outcome. A valid argument against a specified candidate is therefore an acceptable research result.

The present framework takes a conservative position on the measurable physics. It adopts general relativity (GR), with a single Lorentzian metric and specified ordinary matter, as the effective gravitational model. The source proposal is kept separate. This separation allows familiar relativistic calculations to be used without implying that the proposed source already explains their origin.

There are consequently two different questions. One is whether the adopted effective model reproduces the observations that the same GR model reproduces. Conditional on identical matter, parameters, solutions, and measurement rules, it does. The other is whether the source assumptions produce that effective model. No construction establishing the second statement is supplied here. The distinction is substantive: agreement obtained by adopting a theory is not evidence that its equations have been derived from another theory.

This paper establishes the vocabulary and the scope of that distinction. Dynamics states the effective action; Consistency examines its local gravitational degrees of freedom; Relativistic Recovery derives representative measurement relations; Geometry explains the observer-congruence language used to interpret them. Those companion analyses concern GR unless a source-to-target connection is explicitly supplied.

\section{Physical intent and intrinsic ordered motion}
The working ontological hypothesis is that the \Aether{} is a physical four-dimensional substrate and that \Aflow{} is its intrinsic ordered motion. Here \emph{substrate} means a proposed underlying physical reality, not an independently observed material. The intended explanation would connect its organization to the spacetime geometry measured through clocks, distances, and light signals. The smooth manifold specified below supplies a mathematical starting point; it does not establish the physical existence of that substrate.

Observers are conceived as part of the same underlying reality, with access to local relations rather than a view from outside it. ``Experienced space'' names the proposed account of local three-dimensional spatial measurements. ``S-time'' names the intended connection between ordered change and temporal experience. These remain working hypotheses. Neither an embedded spatial slice nor a calibrated clock functional follows merely from giving them names.

The flow is not intended as a wind through an already existing three-dimensional space. Rather, the research asks whether intrinsic source relations can generate the metric, the structure that assigns spacetime intervals and causal directions. This reverses the explanatory order of ordinary motion described using a metric and a clock. A mathematical definition of source motion must therefore state its ordering or parameter structure without silently treating that structure as already measured physical time.

This distinction preserves the physical intent without treating it as an accomplished mechanism. Calling the source four-dimensional does not supply its laws, explain how observers are represented, or establish agreement with GR. Those are reconstruction requirements. The effective model below is the comparison target, not evidence that the proposed source already produces it.

It is also important not to confuse this use of \Aether{} with Einstein-aether theory. The latter introduces a dynamical unit timelike vector with an action and additional coupled modes \parencite{JacobsonMattingly2001,JacobsonMattingly2004}. The effective model here contains no corresponding independent gravitational vector. An observer's four-velocity will appear in the geometric description, but choosing observers is not the same operation as adding a field equation for a new physical substance.

\section{The source assumptions}
\begin{definition}[Source arena]
The symbol $\Sub$ denotes a smooth four-dimensional manifold. Its manifold structure and dimension are assumed. No particular topology is selected by this statement, and no Lorentzian metric, causal order, physical clock, or matter content is assigned to it.
\end{definition}
The source text describes these assumptions as \emph{primitive debt}: they are structure used at the beginning, not structure recovered by a derivation. Here that phrase means only ``assumed rather than explained.'' A later proposal may justify or replace these assumptions, but they cannot currently be advertised as consequences of \Aflow{}.

\begin{definition}[Unspecified source structure]
$\PhiSrc$ is a placeholder for a possible source-order or evolution structure. Its mathematical type has not been fixed. It is not yet a map, vector field, differential operator, stochastic process, group action, or physical-time parameter.
\end{definition}
The distinction between a placeholder and a mathematical flow matters. A conventional flow might be specified by maps $\Phi_\lambda$ obeying a composition law and defined on a parameter domain. Writing such maps here would add assumptions about that domain, composition, and perhaps invertibility. The source definition explicitly leaves those choices unresolved. We therefore retain $\PhiSrc$ without assigning it a convenient but unsupported type.

The source package can be summarized as
\begin{equation}
\Sub\quad\hbox{smooth},\qquad \dim\Sub=4,\qquad
\PhiSrc\quad\hbox{not yet typed}.
\end{equation}
This is a statement of the current setup, not a closed source theory. It supplies no state space for evolving degrees of freedom, no admissible variations, no dynamics, no observables, and no map to clocks or detectors. These omissions constrain what can be inferred; they do not establish that every future specification will fail.

\section{Source arena and physical spacetime}
Let $M$ denote the effective spacetime manifold and $g_{\mu\nu}$ its metric. Writing $(M,g)$ separately from $\Sub$ prevents an ambiguity present in informal descriptions of the ontology. Both arenas may be four-dimensional, but equal dimension does not identify them. Nor does smoothness supply a metric or a causal cone.

A source-to-target reconstruction would need to explain how some independently defined source data determine, approximate, or otherwise represent $(M,g)$ and its operational content. Whether this requires a map between manifolds, a map between state spaces, a quotient, or a limiting construction depends on the source model. No one of these possibilities is adopted here.

It follows that a familiar GR congruence $u^\mu$ cannot simply be declared equal to $\PhiSrc$. The former is a well-defined vector field on an already Lorentzian spacetime; the latter has neither a type nor a target map. The geometric use of the word flow here is therefore an effective description of observer motion, not a construction of source motion.

A further distinction concerns reality claims. Selecting $\Sub$ as a mathematical arena does not establish that a physical substrate with that structure exists. Its proposed physical role is an ontological hypothesis. The effective metric, in contrast, already has an operational role within the adopted model: it enters the rules used to compare clocks, trajectories, and light signals. A future bridge must connect these roles rather than relying on the common use of geometric language.

\section{What ``experienced space'' would need to mean}
The phrase \emph{experienced spatial slice} originally expressed the aim of explaining why observers describe local measurements using three spatial directions. GR already provides a precise local construction for that purpose. Given a future-directed four-velocity normalized by $u^\mu u_\mu=-c^2$, define
\begin{equation}
h_{\mu\nu}=g_{\mu\nu}+\frac{u_\mu u_\nu}{c^2}.
\end{equation}
The projector $h^\mu{}_{\nu}$ selects vectors orthogonal to $u^\mu$. Its image is a three-dimensional local rest space. This is an observer-dependent tangent-space statement, not a new source derivation \parencite{EllisVanElst1999}.

Three different ideas must not be merged. A rest space specifies local spatial directions. A spacelike hypersurface organizes events into a spatial surface when the required integrability conditions hold. A visual scene is reconstructed from received signals and need not correspond to events simultaneous in that rest frame. The ontology has not shown that any of these is an embedded slice of $\Sub$.

For a smooth timelike congruence, vanishing vorticity gives local hypersurface orthogonality. A global slicing needs additional global conditions. The present framework does not assume that every observer family defines a global three-space, and it does not derive a source embedding $\Sigma\hookrightarrow\Sub$. A future account of experienced space must state which of these operational or geometric targets it aims to reconstruct.

\section{S-time: order, duration, and experience}
S-time is retained as the name of an interpretive aim: relating temporal experience to the organization of change. No S-time functional is defined by the present source package. In particular, no monotonicity law or calibration identifies source order with measured duration.

Order and duration are different. An order may distinguish which events can precede others without assigning a number of seconds to an interval. In relativity, proper time assigns duration along a timelike trajectory using the metric. A thermodynamic arrow introduces further questions about irreversible behavior and boundary conditions. Subjective temporal experience raises questions not settled by the mere existence of a clock variable. The source manuscripts contain no derivation connecting all these notions.

The effective model therefore uses ordinary proper time. For a timelike worldline $\Gamma$,
\begin{equation}
\tau[\Gamma]=\frac{1}{c}\int_\Gamma
\sqrt{-g_{\mu\nu}\,\dd x^\mu\dd x^\nu}.
\label{eq:propertime}
\end{equation}
This operational prescription is adopted from GR \parencite{Carroll1997}. Calling it an expression of S-time is an interpretation only until a source construction reproduces the same clock readings and their dependence on trajectory. No universal preferred clock or source time is inferred from the notation.

\section{The adopted relativistic benchmark}
\subsection{Fields, action, and conventions}
The effective arena is a time-oriented four-dimensional Lorentzian spacetime with signature $(-+++)$ and Levi-Civita connection. Greek indices run from $0$ to $3$. In the covariant action, coordinates are length-valued, with $x^0=ct$; $S_m$ has the units of physical action. The benchmark is
\begin{equation}
S_{\mathrm{eff}}[g,\psi]=\frac{c^3}{16\pi G}
\int_M\dd^4x\sqrt{-g}(R-2\Lambda)+S_m[g,\psi],
\label{eq:SA}
\end{equation}
where $G>0$, $\Lambda$ is constant, and $\psi$ denotes specified matter fields. Taking $\Lambda=0$ gives the original noncosmological specialization. Matter is coupled to the same metric, with ordinary locally Lorentz-invariant matter laws when the standard light and clock interpretations are used.

With the physical-action stress tensor
\begin{equation}
T_{\mu\nu}=-\frac{2c}{\sqrt{-g}}\frac{\delta S_m}{\delta g^{\mu\nu}},
\end{equation}
the bulk field equation is
\begin{equation}
G_{\mu\nu}+\Lambda g_{\mu\nu}=\frac{8\pi G}{c^4}T_{\mu\nu}.
\label{eq:einstein}
\end{equation}
Dynamics gives the variation and its boundary assumptions. No source variable occurs in this action as a new propagating gravitational field.

\subsection{Operational content and equivalence}
The effective metric supplies the causal classification of tangent directions. Standard Maxwell light in vacuum geometric optics follows null directions,
\begin{equation}
g_{\mu\nu}k^\mu k^\nu=0,
\label{eq:null}
\end{equation}
and ideal clocks measure \eqref{eq:propertime}. Neutral structureless test bodies follow timelike geodesics when nongravitational forces and finite-size effects can be neglected. These are model assumptions and standard consequences, not new laws induced by an already defined source flow.

The term \emph{exact closure} names the decision to fix this effective gravitational model to GR. For a specified matter model, the two descriptions have identical observables only when they also use the same constants, admissible solutions, initial and boundary data, and measurement rules. If a source interpretation imposed an additional restriction on possible solutions, that restriction would have to enter the comparison. Equality of the printed metric equation alone would no longer suffice.

The action with unspecified $S_m$ thus defines a family of possible GR models. It does not select one cosmological history or one laboratory outcome. Its purpose is to provide a reference against which a completed source model could be compared.

\section{Expansion, light, and the limits of interpretation}
The motivating ontology includes expansion, light propagation, and temporal change. In the adopted benchmark, cosmological expansion is described by a changing FLRW scale factor and the expansion of comoving worldlines. No external container is needed for that description. Light propagation and gravitational redshift follow the same metric and observer measurements developed in Relativistic Recovery.

These facts provide a vocabulary for the proposed interpretation, not evidence that $\PhiSrc$ produces them. A comoving expansion scalar $\theta=3H$, for example, is calculated after a spacetime metric and congruence have been supplied. It cannot serve as the missing source dynamics without reversing the intended explanatory direction.

The same caution applies to metaphor. Static support, rotational observer motion, and wave-induced deformation can all be described using congruences, but no single image of a vortex exhausts gravity. Geometry makes these distinctions quantitative. Its quantities depend on the chosen observers as well as the metric and are not separate manifestations of an independently established source field.

\section{How a later source proposal should be assessed}
A useful source development would first define the objects left open here. It would then state a particular reconstruction claim and show which source assumptions imply it. The claim might concern a causal structure, an operational clock, a metric, or a field equation; these are different levels of achievement. Recovering one does not automatically recover the others.

A successful bridge must avoid importing its intended conclusion as a disguised premise. Starting with a Lorentzian source metric and identifying it with $g$ may define a model, but it would not derive Lorentzian geometry from the present metric-free source assumptions. Similarly, postulating the Einstein--Hilbert action under a different name is not a derivation of that action. Such choices can be legitimate assumptions if they are acknowledged as assumptions.

If an optional extension fails a consistency or observational test, the independent GR benchmark remains a reference. This is a research policy, not a theorem protecting all source hypotheses. A failure may provide evidence against the source assumptions responsible for it. Conversely, absence of a bridge in these manuscripts is not a proof that no bridge can be constructed.

\section{Research purpose and the assumption boundary}
The project seeks to determine whether a mathematically specified development of the ontology can recover GR, or whether a stated reconstruction claim can be ruled out within a defined scope. A negative result is an admissible scientific outcome. The project also provides a proposed research case study for artificial intelligence (AI): can a system formulate precise questions, track its premises, and produce arguments that withstand independent checking? These objectives are related but must be assessed separately.

Sys4AI is the intended agentic research framework. In this role, it would coordinate AI agents, research tasks, records of assumptions, and verification work. That intended use does not supply a physical law or establish the capability of an implemented system. The manuscripts report no completed Sys4AI evaluation runs. The protocol in the synthesis article specifies evidence that such runs would need to produce.

A research task begins with a fixed version of the source assumptions. It must identify the mathematical type of every object used, the admissible source states or histories, and the rule connecting those objects to the proposed target. Smoothness and four-dimensionality remain assumed unless that task explicitly replaces them. Choosing a type for $\PhiSrc$, a topology, an evolution rule, a measure, or a metric is an addition to the present specification, not a deduction from the word flow.

There are two distinct modes of work. In a \emph{fixed-assumption audit}, the agent may analyze the stated premises but may not change them to obtain the desired conclusion. In \emph{open construction}, the agent may propose new definitions and physical assumptions. Each addition must be recorded, motivated, and included in the final theorem statement. Open construction is part of the research objective; it is not permission to describe a result obtained from strengthened premises as a deduction from the original ontology alone.

The first mathematical deliverable should therefore be a source specification and a precise recovery claim, together with a nonempty example or an appropriate existence argument. An inconsistent premise set cannot earn recovery credit through a vacuous implication. Where consistency or existence remains unproved, that dependency stays explicit. A conditional theorem can still be reported as conditional.

\section{Recovery targets and negative conclusions}
A recovery task must name its target before its outcome is assessed. A causal structure describes which events can influence others. A metric assigns spacetime intervals; an operational clock rule relates those intervals to measured duration. Gravitational dynamics specifies how geometry responds to matter. Matter coupling and measurement rules determine which quantities an experiment would observe. Progress on one target does not establish the others.

The task must also identify its domain. A local construction, a special family of solutions, an approximate limiting theory, and exact recovery of a specified GR solution class are different claims. A limited theorem can be valuable without being relabeled as full GR. A derivation of every matter field or the numerical values of every constant is not required unless those are included in the stated claim; externally supplied inputs must nevertheless remain visible.

Four outcomes need separate labels. A \emph{verified derivation} establishes the declared target from the declared premises. A \emph{counterexample to entailment} satisfies the premises but violates the target, showing that those premises do not force that conclusion. A \emph{scoped obstruction} proves that no permitted reconstruction achieves the target within the specified model class and conditions. An \emph{unresolved search} supplies neither the requested derivation nor a valid negative argument.

Non-entailment is weaker than incompatibility. A set of premises may admit a completion that reproduces GR and another that does not. The existence of the latter refutes a universal recovery claim, but not the existence of the former. Conversely, finding one GR-compatible completion does not show that the original ontology selects it uniquely. A proof excluding all admissible completions must quantify over the class actually defined, not over unspecified future meanings of $\PhiSrc$.

These distinctions set the present status. The source specification lacks the definitions needed for the intended reconstruction; this paper does not supply a countermodel to a completed source theory or a theorem excluding every completion. Recognizing that gap is an audit finding. New research progress requires a further checkable result about an explicitly stated development of the proposal.

\section{Conclusion}
The present ontology consists of a selected smooth four-dimensional source arena, an untyped \Aflow{} placeholder, and an interpretive aim concerning spatial and temporal organization. Its quantitative effective model is GR with specified ordinary matter. The geometric and operational results developed in the companion papers belong to that adopted model.

This division preserves the source question without overstating its answer. The existing assumptions do not establish a physical substrate, a source clock, or a derivation of Einsteinian dynamics. They define where further explanation is sought. A future result must supply the missing source structure and demonstrate a specific bridge; until then, exact agreement with GR is a property of the adopted benchmark rather than independent confirmation of the ontology.
The research question remains open in a specific sense: the present source package is not a completed theory from which the target can be decided. The next useful result is a theorem or counterexample about a precisely defined source development. A valid scoped negative finding can advance that inquiry and the assessment of AI reasoning, without being a universal verdict on every possible completion.

\References
\end{document}
